\section{طبقه‌بندی با \lr{KNN}}

در این بخش، الگوریتم \lr{K-Nearest Neighbors (KNN)} برای طبقه‌بندی بیماران بر اساس اندازه‌گیری‌های
بالینی آن‌ها به کار گرفته می‌شود. \lr{KNN} یک روش یادگیری مبتنی بر نمونه و ناپارامتری است و برای این
مجموعه‌داده مناسب است؛ زیرا در تحلیل اکتشافی داده‌ها تفکیک‌پذیری خطی قوی میان کلاس‌ها مشاهده نشد و
چندین ویژگی، به‌ویژه \lr{thalach} و \lr{oldpeak}، روابط غیرخطی با وجود بیماری قلبی نشان می‌دهند.
از آن‌جا که \lr{KNN} مستقیماً بر فاصله بین بردارهای ویژگی تکیه دارد، پیش‌پردازش مناسب و مقیاس‌بندی
درست ویژگی‌ها اجزای کلیدی زنجیرهٔ طبقه‌بندی محسوب می‌شوند.

\subsection{مروری بر الگوریتم \lr{KNN}}

برای یک نمونهٔ آزمون $x$، طبقه‌بند \lr{KNN} به صورت زیر عمل می‌کند:

\begin{enumerate}
	\item محاسبهٔ فاصله بین $x$ و تمامی نمونه‌های مجموعهٔ آموزش،
	\item انتخاب $k$ نمونهٔ آموزشیِ نزدیک‌تر به $x$،
	\item تعیین کلاس غالب در میان این $k$ همسایه،
	\item انتساب همان برچسب کلاس به نمونهٔ $x$.
\end{enumerate}

از آن‌جا که \lr{KNN} هیچ فرض مشخصی دربارهٔ شکل تابعی یا توزیع احتمالاتی داده‌ها ندارد، برای
مجموعه‌داده‌های پزشکی که در آن‌ها برهم‌کنش‌های پیچیده و غیرخطی میان متغیرها وجود دارد، بسیار
مناسب است؛ موضوعی که در نمودارهای جفتی و الگوهای همبستگی مشاهده‌شده در بخش تحلیل اکتشافی نیز
تأیید شد.

\subsection{تابع فاصلهٔ ترکیبی برای ویژگی‌های عددی و دسته‌ای}

مجموعه‌دادهٔ مورد استفاده شامل ترکیبی از ویژگی‌های عددی (مانند \lr{age}، \lr{trestbps}،
\lr{chol}، \lr{thalach} و \lr{oldpeak}) و ویژگی‌های دسته‌ای یا دودویی (مانند \lr{sex}،
\lr{cp}، \lr{exang}، \lr{restecg}) است. استفاده از یک معیار فاصلهٔ صرفاً عددی (مثلاً
فاصلهٔ اقلیدسی روی همهٔ ستون‌ها) به‌تنهایی مناسب نیست و می‌تواند به تفسیر نادرست
ویژگی‌های دسته‌ای منجر شود.

برای رفع این مشکل، یک تابع فاصلهٔ ترکیبی به کار گرفته شد که در آن برای ویژگی‌های عددی از فاصلهٔ
اقلیدسی و برای ویژگی‌های دسته‌ای از فاصلهٔ \lr{Hamming} استفاده می‌شود:

\[
d(x,y) = w_{\text{num}} \cdot d_{\text{num}}(x,y) +
w_{\text{cat}} \cdot d_{\text{cat}}(x,y),
\]

که در آن، هر دو وزن در این مطالعه برابر $0.5$ در نظر گرفته شدند. این وزن‌دهی باعث می‌شود ویژگی‌های
عددی—به‌خصوص ویژگی‌های تأثیرگذاری مانند \lr{thalach} و \lr{oldpeak}—در محاسبهٔ شباهت نقش پررنگی
داشته باشند، بدون آن‌که اطلاعات موجود در ویژگی‌های دسته‌ای نادیده گرفته شود.

\subsection{آموزش مدل و انتخاب مقدار $k$}

پس از انجام پیش‌پردازش و نرمال‌سازی \lr{Min--Max}، مجموعه‌داده با نسبت ۸۰/۲۰ به آموزش و آزمون
تقسیم شد. طبقه‌بند \lr{KNN} روی ماتریس \lr{X\textsubscript{train}} آموزش داده شد و روی
\lr{X\textsubscript{test}} مورد ارزیابی قرار گرفت.

برای بررسی تأثیر اندازهٔ همسایگی، الگوریتم \lr{KNN} برای مقادیر مختلف $k$ ارزیابی شد:

\[
k \in \{3,\,5,\,7\}.
\]

برای هر مقدار $k$، عملکرد مدل با استفاده از معیارهای دقت (\lr{Accuracy})، دقت مثبت (\lr{Precision})،
بازخوانی (\lr{Recall}) و امتیاز \lr{F1} سنجیده شد؛ معیارهایی که در کاربردهای بالینی اهمیت ویژه‌ای دارند،
زیرا هم خطاهای نوع اول (تشخیص اشتباه بیماری) و هم خطاهای نوع دوم (عدم تشخیص بیماری) می‌توانند
پیامدهای جدی داشته باشند.

\subsection{نتایج طبقه‌بندی}

\begin{table}[H]
	\centering
	\caption{عملکرد طبقه‌بند \lr{KNN} برای مقادیر مختلف $k$.}
	\label{table:knn-results}
	\begin{tabular}{cccccc}
		\toprule
		$k$ & Accuracy & Precision & Recall & F1-score \\
		\midrule
		3 & 0.9317 & 0.9159 & 0.9515 & 0.9333 \\
		5 & 0.8195 & 0.7946 & 0.8641 & 0.8279 \\
		7 & 0.8341 & 0.8053 & 0.8835 & 0.8426 \\
		\bottomrule
	\end{tabular}
\end{table}

بیشترین مقدار دقت در $k = 3$ به دست آمد. این نتیجه با سطح متوسط هم‌پوشانی کلاس‌ها که در نمودارهای
پراکندگی و نمودارهای جفتی مشاهده شد سازگار است؛ مقادیر کوچک‌تر $k$ به طبقه‌بند اجازه می‌دهند الگوهای
محلی را که عمدتاً توسط دو ویژگی بسیار اطلاع‌بخش \lr{thalach} و \lr{oldpeak} شکل می‌گیرند، بهتر
شناسایی کند. با افزایش $k$، این تمایزهای محلی بیش از حد هموار شده و کارایی مدل کاهش می‌یابد.

در نتیجه، اندازهٔ همسایگی بهینه برای این مجموعه‌داده برابر است با:

\[
k_{\text{best}} = 3.
\]

\subsection{ماتریس اغتشاش}

ماتریس اغتشاش برای بهترین مدل (با $k = 3$) در شکل \ref{fig:confusion-best} نشان داده شده است.

\begin{figure}[H]
	\centering
	\includegraphics[width=0.6\textwidth]{confusion_matrix.png}
	\caption{ماتریس اغتشاش طبقه‌بند \lr{KNN} برای $k=3$.}
	\label{fig:confusion-best}
\end{figure}

همان‌طور که مشاهده می‌شود، مدل تعادلی مناسب بین حساسیت و ویژگی (\lr{Sensitivity/Specificity}) برقرار
کرده است. مقدار نسبتاً بالای بازخوانی نشان می‌دهد که بیشتر موارد بیماری قلبی به‌درستی شناسایی شده‌اند.
در عین حال، تعداد کمی خطای نوع اول (مثبت کاذب) و نوع دوم (منفی کاذب) باقی می‌ماند که با توجه به
هم‌پوشانی توزیع‌ها در تحلیل اکتشافی، به‌ویژه در ویژگی‌هایی مانند \lr{age} و \lr{chol} که جداسازی
ضعیف‌تری دارند، قابل انتظار است. با وجود این، مدل در مجموع عملکرد پیش‌بینی‌کنندهٔ قدرتمندی را با توجه
به ماهیت ترکیبی داده‌ها و تعداد محدود ویژگی‌های واقعاً تفکیک‌کننده نشان می‌دهد.